\section{KVA、页表回收与动态内核栈}

\topic{not-present 从来不等于 free}
早期内核的 heap、初始栈、IST、LAPIC 和写保护测试页都有固定地址。只要每种
资源恰好存在一份，链接脚本和常量表就能避免冲突；动态 Thread、可增长 slab、
临时设备缓冲和模块映射出现后，地址本身也成为必须申请、持有和归还的资源。

页表不能承担这份软件所有权。PTE 的 present 位只回答处理器现在能否把线性
地址翻译成物理地址。以下三种状态在硬件看来都可能是 present=0，软件含义却
完全不同：

\begin{osgridlongtable}{L{0.24\textwidth}L{0.22\textwidth}L{0.40\textwidth}}
\textbf{软件状态} & \textbf{叶 PTE} & \textbf{能否交给新调用者}\\ \hline
\endfirsthead
\textbf{软件状态} & \textbf{叶 PTE} & \textbf{能否交给新调用者}\\ \hline
\endhead
空闲地址 & not-present & 可以，由 KVA 分配器提交所有权\\ \hline
事务准备中 & not-present & 不可以，调用者正在取得物理后备或构造页表\\ \hline
guard page & 永久 not-present & 不可以，它属于栈区间并负责捕获越界\\ \hline
\end{osgridlongtable}

若通过遍历 PTE 寻找 not-present 空洞，两个并发或交错事务都可能在首个映射
提交前得到同一地址；guard 也会被错误复用。KVA（kernel virtual address）
分配器因此只回答一个问题：“哪段虚拟页区间已经许诺给谁？”它不申请物理页、
不设置 present，也不选择 RW/NX 权限。

这种分层不是本项目发明的偶然术语。早期 UNIX 内核主要依赖静态映射；随着
内核模块、网络缓冲和动态对象增加，内核需要把不连续物理页组织成连续虚拟
区间。现代系统常把这类能力归入 vmalloc 一族，并加入树索引、延迟 TLB 刷新、
调试 guard 与并发回收。本项目先实现串行、有界的区间状态机，
不复制尚无调用场景的性能层。

\topic{32 TiB 窗口怎样落在四级页表中}
四级分页使用 48 位 canonical 线性地址：bit 63..48 必须全部复制 bit 47。
高半区因此从 \code{0xFFFF800000000000} 延伸到
\code{0xFFFFFFFFFFFFFFFF}。当前内核把 direct-map 与 KVA 排列为：

\begin{osgridlongtable}{L{0.20\textwidth}L{0.24\textwidth}L{0.24\textwidth}L{0.18\textwidth}}
\textbf{区域} & \textbf{起点} & \textbf{末端（不含）} & \textbf{PML4 槽}\\ \hline
\endfirsthead
\textbf{区域} & \textbf{起点} & \textbf{末端（不含）} & \textbf{PML4 槽}\\ \hline
\endhead
direct-map &
\code{FFFF888000000000} &
\code{FFFFC88000000000} &
273--400\\ \hline
KVA &
\code{FFFFC90000000000} &
\code{FFFFE90000000000} &
402--465\\ \hline
\end{osgridlongtable}

一个 PML4 槽覆盖 512 GiB，64 个槽正好覆盖 32 TiB。KVA 起点的四级索引是
\((402,0,0,0)\)，exclusive 末端是 \((466,0,0,0)\)。起点、长度和末端都在
提交前按 64 位无符号数验证：页数乘 4096 不得溢出，末端加法不得回绕，末页
仍必须满足 48 位 canonical 规则。未来若启用五级分页 LA57，需要新的布局
迁移，不能让同一个 ABI 常量在两种符号扩展规则下静默改变含义。

窗口包含
\[
 \frac{32\ \mathrm{TiB}}{4\ \mathrm{KiB}}
 =8589934592=2^{33}
\]
个虚拟页。首个 4 KiB 页永久标记为 reservation。它不占物理页，也没有叶
映射；它只让窗口零偏移误用具有稳定的软件诊断，而不是直接成为第一个正常
对象。

\topic{为什么 32 TiB 没有配一张逐页位图}
若每个虚拟页只用一位表示“占用”，元数据也需要
\[
 \frac{2^{33}}{8}=2^{30}\ \mathrm{bytes}=1\ \mathrm{GiB}.
\]
要区分 reservation 与 allocation 至少还要增加状态位。启动阶段实际同时
存在的 KVA 对象远少于数十亿页，为稀疏可能空间支付 GiB 常驻元数据没有比例
意义。

当前实现由调用者提供 256 项
\code{KernelVirtualAddressRangeDescriptor} BSS 数组。每项恰有三个 64 位
字段：

\begin{center}
\begin{osgridtabular}{lll}
字段 & 含义 & 有效约束\\ \hline
\code{begin\_address} & 区间首虚拟地址 & 4 KiB 对齐、位于窗口内\\ \hline
\code{page\_count} & 区间页数 & 非零、换算字节不溢出\\ \hline
\code{kind} & Unused/Allocation/Reservation & 活动前缀不能为 Unused\\ \hline
\end{osgridtabular}
\end{center}

总大小为 \(256\times24=6144\) 字节。数组前部是按起始地址严格递增的活动
描述符，后部必须保持零地址、零页数和 Unused。256 限制的是“同时存在多少
区间”，不限制单个区间能有多少页。描述符耗尽返回
\code{MetadataExhausted}；窗口没有足够连续页返回
\code{OutOfVirtualAddressSpace}。区分这两种失败，才能判断应扩元数据还是
回收地址。

空闲区间不保存节点，而由窗口边界和相邻有主区间的缝隙推导：

\begin{center}
\begin{tikzpicture}[os diagram,node distance=0mm]
  \node[os state,minimum width=2.2cm] (r) {reservation};
  \node[os block,right=of r,minimum width=2.5cm] (f1) {free gap};
  \node[os state,right=of f1,minimum width=2.2cm] (a) {allocation};
  \node[os block,right=of a,minimum width=2.5cm] (f2) {free gap};
\end{tikzpicture}
\end{center}

删除一个 allocation 描述符后，两侧 gap 在模型上立即连成一个空闲区间，
不存在另一条空闲链需要合并。代价是申请要扫描至多 256 项，插入和删除要移动
尾部，复杂度为 \(O(N)\)。当前串行规模下，完整可检查的不变量比尚未被测量
需要的红黑树更重要。

\topic{best-fit 必须对齐绝对虚拟页号}
请求页数 \(n\)、对齐页数 \(a\) 时，\(a\) 必须是二的幂。对一个字节起点为
\(g\) 的 free gap，先得到绝对虚拟页号
\[
 p=\frac{g}{4096},
 \qquad
 p_{\mathrm{aligned}}=(p+a-1)\mathbin{\&}\sim(a-1),
\]
再把它安全乘回字节地址。这里不能对
\((g-\mathrm{window\_begin})/4096\) 做相对对齐；“8 页对齐”的契约是返回
地址本身能被 32 KiB 整除，不是它与窗口起点的距离恰好整除。

每个 gap 都先检查对齐前缀、剩余页数和溢出。所有可容纳候选中选择原始 gap
页数最小者；同样大小时保留扫描得到的低地址。这是 best-fit：优先消耗最小
可用洞，减少把大洞用于小请求的概率。它不是消除碎片的数学保证，但对固定
上限提供确定、可由独立模型复现的结果。

候选地址、插入位置和计数增量全部验证后，实现才移动描述符并填写调用者输出。
任何失败都保持数组、统计和输出 range 原值。这样上层跨资源事务可以根据
明确状态执行回滚，不会拿到“返回失败但部分有效”的地址。

\topic{释放 KVA 时为什么必须给出原区间}
\code{ReserveRange} 用来表达固定软件占用。它与普通申请共享页对齐、窗口、
canonical、溢出和重叠检查，但不能由 \code{TryRelease} 释放。重叠检查先于
元数据容量检查，因此即使 256 项已满，重复保留已有地址也仍报告
\code{RangeOverlap}，而不是被较表面的容量错误遮蔽。

释放必须同时匹配原描述符的起始地址、页数和 Allocation 类型。区间内部地址
返回 \code{AllocationNotFound}，正确起点但错误页数返回
\code{AllocationSizeMismatch}，reservation 返回 \code{ReservedRange}，
已经删除的地址再次释放仍返回 \code{AllocationNotFound}。实现不会向前寻找
“看起来像块头”的数据，也不会猜测调用者希望释放整个包含区间。

成功删除后，数组尾部向前压紧，最后一项恢复全零。累计统计始终满足
\[
 \mathrm{successful\ allocations}
 -
 \mathrm{releases}
 =
 \mathrm{active\ allocations}.
\]
当前分配页与保留页之和不得超过 \(2^{33}\)，当前值不得超过各自峰值。

\topic{Validate 从唯一有主集合重建全部事实}
\code{Validate} 不申请临时内存。它先验证窗口本身，再沿活动前缀逐项证明：

\begin{enumerate}[label=\arabic*.]
  \item 类型只能是 Allocation 或 Reservation，页数非零，区间完整位于窗口；
  \item 当前起点不小于前一区间末端，因此允许紧贴但禁止倒序和重叠；
  \item 重新累计分配页、保留页、活动申请和 reservation 数；
  \item 活动前缀之后每个描述符都严格为 Unused、零地址、零页数；
  \item 从窗口首尾与相邻描述符重新计算 free page 和最大连续 gap；
  \item 把重建结果与缓存统计、累计申请/释放和峰值逐一比较。
\end{enumerate}

空闲 gap 没有第二份持久索引，所以不会出现“活动树正确、空闲链漏了一块”的
双结构漂移。若未来为性能加入平衡树或分段索引，它必须被视为可重建的派生
结构，完整验证仍应回到唯一所有权集合。

\topic{映射失败时怎样撤销新页表}
取得 KVA 只是跨层事务的第一步。一个带 guard 的动态对象遵循：

\begin{verbatim}
allocate KVA range
  -> allocate physical block or pages
  -> map data pages with explicit permissions
  -> publish object
\end{verbatim}

任何中间失败都要从已提交的最后一步开始撤销。正常销毁同样固定为：

\begin{verbatim}
unpublish object
  -> clear leaf PTEs and invalidate TLB entries
  -> return physical pages to buddy
  -> release exact KVA range
\end{verbatim}

顺序不能交换。若先归还物理页，旧 TLB 翻译仍可能写入已经交给新对象的页；
若先释放 KVA，另一个调用者可以在旧映射尚存时取得同一虚拟地址。guard page
从申请到释放始终没有叶映射，却与数据页一起属于同一个 KVA 区间，这正是
“not-present 不等于 free”的执行例子。

当前 \code{UnmapPage} 不只清除叶项。它先证明整条路径有效且属于当前页表根，
再按 PT、PD、PDPT 从子到父回收已经变空的独占分支。KVA 自检先用一个虚拟页
和一个 order-0 物理页暖机首条页表路径：撤销时回收 PT 与 PD，却保留可能仍
被进程 PML4 项引用的共享 PDPT。主事务随后复用该共享入口，重新创建并在最后
一张叶撤销时回收 PT 与 PD。因此物理页、buddy、KVA 与页表自身都进入可解释
的所有权基线；唯一保留的 PDPT 是共享协议，不是被忽略的泄漏。

\topic{用六页事务同时检查 guard、权限和回收}
预热后，QEMU 中的正式 KVA 自检执行：

\begin{enumerate}[label=\arabic*.]
  \item 申请 6 页、8 页对齐区间，得到
        \code{0xFFFFC90000008000}；
  \item 从 buddy 申请 order 2，即 4 个连续物理页；
  \item 相对页 0 和 5 保持 not-present，作为两个 guard；
  \item 把相对页 1--4 映射为 supervisor、RW、NX；
  \item 用 \code{QueryPage} 核对四个物理地址连续、页粒度和权限精确；
  \item 经首尾数据虚拟页写入两个 64 位模式并读回；
  \item 逆序撤销四个映射，归还整个 order-2 块，再精确释放六页 KVA；
  \item 比较三层基线，运行 buddy 与 KVA 的完整校验，并核对两段事务累计
        回收两张 PT、两张 PD、零张 PDPT，精确保留一张共享 PDPT。
\end{enumerate}

只有事务全部提交，冷路径才输出一次稳定快照：

\begin{verbatim}
[OS][KERNEL] KVA_ALLOCATOR_READY
[OS][KERNEL] KVA_WINDOW_BASE=0xFFFFC90000000000
[OS][KERNEL] KVA_WINDOW_SIZE_BYTES=0x0000200000000000
[OS][KERNEL] KVA_DESCRIPTOR_CAPACITY=0x0000000000000400
[OS][KERNEL] KVA_ACTIVE_DESCRIPTORS=0x0000000000000001
[OS][KERNEL] KVA_ALLOCATED_PAGES=0x0000000000000000
[OS][KERNEL] KVA_RESERVED_PAGES=0x0000000000000001
[OS][KERNEL] KVA_SUCCESSFUL_ALLOCATIONS=0x0000000000000002
[OS][KERNEL] KVA_RELEASES=0x0000000000000002
[OS][KERNEL] KVA_PEAK_ALLOCATED_PAGES=0x0000000000000006
[OS][KERNEL] KVA_SELF_TEST_MAPPED_PAGES=0x0000000000000004
[OS][KERNEL] KVA_SELF_TEST_GUARD_PAGES=0x0000000000000002
[OS][KERNEL] KVA_SELF_TEST_PASSED
[OS][KERNEL] PAGE_TABLE_RECLAIMED_LEVEL1_TABLES=0x0000000000000002
[OS][KERNEL] PAGE_TABLE_RECLAIMED_LEVEL2_TABLES=0x0000000000000002
[OS][KERNEL] PAGE_TABLE_RECLAIMED_LEVEL3_TABLES=0x0000000000000000
[OS][KERNEL] PAGE_TABLE_RETAINED_SHARED_LEVEL3_TABLES=0x0000000000000001
[OS][KERNEL] PAGE_TABLE_RECLAIM_SELF_TEST_PASSED
\end{verbatim}

申请、描述符移动、映射和释放热路径不逐步打串口。日志只证明最终事务和守恒，
不会让 115200 波特 I/O 改变分配器时序。

\topic{三类宿主模型与真实处理器路径各自证明什么}
单元测试覆盖未初始化、无存储、零容量、canonical 空洞、保留重叠、best-fit、
绝对对齐、输出失败原子性、描述符与地址分别耗尽、精确释放错误和人为破坏
元数据。它证明边界和每个状态码，但不宣称页表真正执行。

集成测试串联真实 \code{PhysicalFrameAllocator}、
\code{KernelVirtualAddressAllocator} 与 \code{PageTableManager}，在宿主
模拟物理内存上执行同一六页生命周期。它核对 guard、四个物理身份、RW/NX、
写入模式和最终三层统计，证明模块组合没有遗漏回滚。

固定种子 \code{0x4B564152414E444F} 在 512 页独立逐页模型上执行 100000 步
申请与释放。参考模型自己扫描 gap、计算 best-fit 与对齐，不调用被测分配器；
每 257 步比较活动集合、累计统计、最大空洞并运行完整校验，结束时再随机顺序
排空所有申请。

QEMU 最后让真实 x86-64 MMU、TLB 和 load/store 执行同一契约。64 MiB 证明
低资源兼容路径，64 GiB 证明现代规格没有改变地址层语义。三类宿主测试加一条
处理器执行路径不是重复劳动：它们分别定位算法、组合、长序列不变量和硬件
效果。

当前 KVA 最多保存 1024 个有主区间，操作为 \(O(N)\)，
没有内部锁。初始上界曾是 256；扩容后可以同时容纳 512 个动态
Thread 栈及其他启动事务，不改变区间算法。
空中间页表回收已经由独立的页表所有权协议闭环；动态内核栈则成为第一个长期
持有 KVA、物理帧和叶映射的真实调度资源。性能树、per-CPU KVA cache 和延迟
shootdown 必须等并发契约与测量证据出现后再加入。

\section{不同页表根怎样共享内核半区}
80386 的 32 位分页已经采用目录加页表的层次结构，避免为整个线性地址空间
预先保存逐页表项。PAE 把表项扩为 64 位并增加层级，x86-64 又让常见四级
模式用四组 9 位索引覆盖 48 位 canonical 线性地址。这个历史方向始终在解决
“怎样让硬件快速找到翻译”，没有替操作系统回答“最后一个映射消失后谁可以
释放中间表”。

在四级 4 KiB 模式中，每张表恰有 512 个 8 字节项，所以正好占一个 4096
字节物理页。各层覆盖范围为：

\begin{osgridlongtable}{@{}lll@{}}
层 & 一个表项覆盖 & 一张完整表覆盖\\ \hline
PT / level 1 & \(4\ \mathrm{KiB}\) & \(2\ \mathrm{MiB}\)\\ \hline
PD / level 2 & \(2\ \mathrm{MiB}\) & \(1\ \mathrm{GiB}\)\\ \hline
PDPT / level 3 & \(1\ \mathrm{GiB}\) & \(512\ \mathrm{GiB}\)\\ \hline
PML4 / level 4 & \(512\ \mathrm{GiB}\) & \(256\ \mathrm{TiB}\) canonical envelope\\ \hline
\end{osgridlongtable}

CR3 给出根表物理地址。处理器从虚拟地址取出 \(47{:}39\)、\(38{:}30\)、
\(29{:}21\)、\(20{:}12\) 四组索引，最后把 \(11{:}0\) 当页内偏移。每个
非叶项主要告诉硬件下一张表的物理页和 P/RW/U/S 等遍历权限；叶项再决定
最终页的 NX、缓存、脏位和访问位等属性。这里没有父指针、引用计数、
\code{owner} 或“这张表由哪个地址空间申请”的字段。

这一缺失不是 x86 设计漏洞。硬件只需要在每次 load/store 时完成确定遍历；
对象生命周期属于操作系统策略。代价是软件不能从某个非叶项孤立地推出子表
是否独占。

\topic{复制 PML4 项共享的是子树，不是子树副本}
当前进程地址空间为了共享内核映射，会把内核根中的高半 PML4 项按值复制到
新的进程 PML4。复制的 64 位值仍包含同一 PDPT 物理地址：

\begin{verbatim}
kernel PML4E ----+
                 +----> shared PDPT ---> PD ---> PT
process A PML4E -+
process B PML4E -+
\end{verbatim}

因此“内核根里这条分支已经没有叶映射”和“没有任何根引用该 PDPT”是两个
不同命题。若内核只清自己的 PML4E 并归还 PDPT，进程 A/B 的 PML4E 仍指向
原物理帧。buddy 可以马上把该帧交给堆或文件缓存；下一次在 A 中发生硬件
page walk 时，处理器会把新对象字节解释为页表项。这是页表帧
use-after-free，通常比普通悬空指针更难定位，因为损坏在另一个 CR3 被调度后
才显现。

相反，PDPT 内某个 PDPTE 位于真正共享的物理页中。清除它会被所有引用该
PDPT 的根共同观察到，所以它指向的空 PD 以及更低的 PT 可以安全释放。这个
差异给出本阶段的核心原则：

\begin{quote}
空只说明没有有效内容；回收还必须证明当前管理器拥有最后一个硬件引用。
\end{quote}

\topic{三种页表根为什么不能混用}
\code{PageTableManager} 构造时必须显式接收
\code{PageTableRootKind}，不能从虚拟地址、当前 CR3 或调用栈猜测：

\begin{osgridlongtable}{@{}p{0.18\textwidth}p{0.42\textwidth}p{0.30\textwidth}@{}}
根类型 & 允许修改的分支 & 最深回收边界\\ \hline
\code{Exclusive}
& 整棵树都只属于该管理器
& PT、PD、PDPT\\ \hline
\code{KernelShared}
& 内核根中的映射；下级父项变化对所有借用者可见
& PT、PD；PDPT 保持稳定\\ \hline
\code{Process}
& PML4[0] 下的用户程序分支和 PML4[255] 下的用户栈分支
& 程序低端克隆 PDPT 留到进程销毁；用户栈可回收到 PDPT\\ \hline
\end{osgridlongtable}

进程低地址还有一个细节。PML4[0] 指向克隆的 PDPT，其中只有约定的
PDPT[1] 用户程序分支归进程修改；其余低端内容可能来自启动兼容映射。用户栈
放在 PML4[255] 的独占分支。对复制来的内核高半分支或其他借用分支调用
\code{MapPage}/\code{UnmapPage}，必须在申请帧和写表项之前返回
\code{SharedBranchMutationDenied}。

把根类型设为构造必选参数还有工程价值：新增调用点若没有回答所有权问题就
无法编译。相比默认成“最宽松”或根据地址猜测，这会把架构判断前移到代码
评审和单元测试。

\topic{“空”必须表示其余 511 个原始项逐字为零}
一个表是否可回收，不能只数 Present 位。现代操作系统经常在 non-present
PTE 中保存 swap 类型、文件页状态或诊断标记；处理器不继续遍历，不代表软件
没有保存事实。当前项目尚未实现 swap，也仍采用更强定义：排除马上要解除的
父子边后，其余每个 64 位项必须等于零。

对普通 4 KiB 叶，判定链为：

\begin{verbatim}
PT except target PTE is all-zero
  -> PD except target PDE is all-zero
    -> root policy permits PDPT reclamation
      and PDPT except target PDPTE is all-zero
\end{verbatim}

一旦某层不空，更高层必然仍有有效子树，检查即可停止。2 MiB 页把 PDE 本身
作为叶并设置 PS；其数据后备通常对应 order 9 连续页。当前
\code{UnmapPage} 有意拒绝这种情况并返回
\code{UnexpectedLargePage}，不把不同粒度的数据所有权偷偷塞进 4 KiB
接口。

\topic{撤销先完整预检，再一次性断开}
直接“边走边清”会产生难以回滚的半状态。假设已经清掉 PTE，走到 PD 时才
发现父项指向未拥有或越界物理帧；调用方既失去映射，又无法证明中间树安全。
当前撤销因此分为两阶段。

只读预检逐层完成：

\begin{enumerate}[label=\arabic*.]
  \item 检查虚拟地址 canonical 且按 4 KiB 对齐；
  \item 检查每个子表物理地址按页对齐，完整落在允许访问范围；
  \item 拒绝子表等于根或任一祖先的明显环；
  \item 用 buddy 的精确查询证明每张子表仍是活动 order-0 allocation；
  \item 保存四级父项地址和三个子表物理身份；
  \item 按原始 64 位值计算 PT、PD、PDPT 是否会在本次撤销后变空；
  \item 在修改任何表项前，再次预检所有计划释放的帧。
\end{enumerate}

预检失败返回 \code{InvalidTableFrame} 或
\code{TableFrameNotOwned}，叶项、调用方结果对象、TLB 与分配器统计全部
保持原样。预检成功后的提交顺序是：

\begin{verbatim}
clear PTE
clear PDE/PDPTE/PML4E that lead only to reclaimable empty tables
INVLPG target virtual address exactly once
zero and release PT, then PD, then permitted PDPT
\end{verbatim}

父项先断开，确保新的硬件遍历无法进入即将归还的子表；页帧再从子到父释放，
与创建依赖顺序相反。一次目标叶修改只需一次
\code{INVLPG}，为每个已释放中间表重复失效同一线性地址没有额外语义。
当前是单 BSP，所以还没有向其他 CPU 发送 TLB shootdown。

\code{PageTableUnmapResult} 分别返回 level 1、level 2、level 3 和总回收
表帧数。数据页由上层对象拥有，不由撤销接口释放，也不计入结果。这使调用者
可以精确区分“PTE 已清但共享表仍存在”和“本次真正回收了页表帧”。

\topic{映射也必须能撤销父项权限与新表}
映射一张原先不存在的叶页，可能依次成功创建 PDPT、PD，却在申请 PT 时耗尽。
旧式实现若立即链接每张新表，只返回
\code{FrameAllocationFailed}，就会永久留下两个空页。另一个更隐蔽的例子是
用户映射：走过既有父项时需要把 U/S 提升为 1；若最后发现叶已经映射，临时
提升也不能残留。

当前每次非叶穿越记录一个 \code{TableMutation}：

\begin{verbatim}
entry                   address of the parent entry
original_entry          complete 64-bit value before mutation
table_physical_address  child table frame
table_created           whether this transaction allocated it
\end{verbatim}

只有叶项写入成功才提交事务。此前遇到帧耗尽、大页冲突、重复映射、损坏地址
或所有权错误时，先反序恢复每个父项完整原值，再反序清零并释放本事务创建的
表帧。保存完整值而非只保存 U/S 位，是为了不假设未来不会加入新的软件位。
若回滚自身遇到帧所有权或释放故障，返回
\code{RollbackFailed}，不能继续用最初错误掩盖已经不可证明的树状态。

这种局部日志与数据库 write-ahead log 的目的相似：事务先保存足以恢复的
旧事实，再改变共享结构。区别是它使用固定长度栈上记录，不分配内存，也没有
持久化恢复需求。

\topic{查询和进程整体销毁也必须服从同一边界}
\code{QueryPage} 是只读接口，却不能无条件解引用页表地址。损坏或已释放的
父项同样会让宿主测试越界、让目标内核读取任意 direct-map 内容。因此查询在
每一级复用范围、对齐、环和精确 order-0 所有权验证。命中 4 KiB 页时返回
\[
  \mathrm{physical\ page\ base} + (\mathrm{virtual\ address}\bmod4096),
\]
不能像旧实现那样遗漏页内偏移；2 MiB 叶使用对应的大页偏移。

逐页撤销是正常细粒度路径，\code{ReleaseProcessRoot} 仍作为地址空间最终
销毁的兜底。它只遍历用户程序与用户栈的自有树，不进入借用的内核高半子树；
已被逐页撤销清空的分支自然跳过，尚存的用户叶和表由递归路径归还。每释放
一张子表都立即清父项，避免部分失败留下指向已归还帧的硬件指针。递归下降
同时携带从根到当前表的完整祖先物理地址链；若下级项回指根或任一祖先，
必须在释放任何相关表帧前返回 \code{InvalidTableFrame}。测试随后恢复原项并
重试销毁，用同一个对象证明拒绝路径没有偷改所有权状态。

\topic{从一次映射到十万步随机序列}
页表回收新增三类宿主测试。单元测试分别建立三种根，覆盖重复初始化、单叶
完整级联、同 PT 相邻叶、共享 PDPT 保留、进程借用拒绝、程序与栈分支差异、
申请耗尽回滚、所有权破坏与原址修复、越界表地址、查询祖先环、递归销毁
祖先回指和失败输出保持。

集成测试不在每轮重建世界，而是在同一分配器上连续执行 128 次共享根生命周期
和 64 次进程根创建、映射、逐页撤销与销毁。这样一个每轮泄漏一页的错误无法
被 fixture 清零隐藏；最后还让进程递归销毁一棵未逐叶排空的树。

固定种子 \code{0x5047545245434C4D} 在 1024 个虚拟页上执行 100000 步。
地址集合跨 4 个 PML4 项、每项 2 个 PDPT 项、多个 PD 与 PT；独立模型只保存
活动叶集合，并由层级索引重算应存在的 PT/PD/PDPT 数，不调用被测回收逻辑。
每 257 步逐页查询、比较表数和 buddy 统计，结束时排空全部映射。

QEMU 最后用真实 CR3、page walk、TLB 和访存验证暖机与四叶 KVA 事务。稳定
摘要为：

\begin{verbatim}
PAGE_TABLE_RECLAIMED_LEVEL1_TABLES = 2
PAGE_TABLE_RECLAIMED_LEVEL2_TABLES = 2
PAGE_TABLE_RECLAIMED_LEVEL3_TABLES = 0
PAGE_TABLE_RETAINED_SHARED_LEVEL3_TABLES = 1
\end{verbatim}

PT 与 PD 各两张来自暖机和主事务；PDPT 为零并精确保留一张，是共享根策略的
正面证据。映射、查询、撤销、空表扫描和随机测试都不逐项写串口，只在完整
初始化提交后输出一次摘要。至此 KVA 与动态栈不再靠暖机泄漏建立基线。

当前边界也必须写清：回收仅支持普通 4 KiB 叶；页帧元数据只证明精确
order-0 allocation，还没有 PageTable 类型标签；\code{KernelShared}
保守保留全部 PDPT，没有通用引用计数；没有 PCID、批量 unmap、远端
shootdown、页表锁或 RCU 读者。这些属于 Process/Thread 与并发地址空间模型
之后的独立问题，不应削弱本阶段已经闭合的单核所有权契约。

\section{动态内核栈怎样跨越特权切换}
x86 的任务状态段诞生于硬件任务切换时代。现代 64 位内核通常不使用 TSS
保存整份任务现场，但处理器从低特权级进入高特权级时，仍需要一个不受用户
RSP 控制的可信落脚点。当前系统在 CPL3 收到中断、异常或
\code{INT 0x80} 时，处理器从本 CPU 的 TSS.RSP0 取 Ring 0 栈顶，再压入
旧 SS、RSP、RFLAGS、CS 和 RIP。汇编入口继续在同一栈上保存通用寄存器。

这形成一个重要的生命周期悖论：进程终止代码正在使用的资源，恰好也是该进程
应当释放的资源。若退出处理函数直接清除当前栈的 PTE，下一条函数返回、局部
变量访问或寄存器恢复就会触发内核页故障；若先把物理页归还 buddy，旧 TLB
翻译还可能写入已经交给其他对象的帧。正确答案不是“销毁函数更小”，而是把
终止分成逻辑终止和物理回收两个时刻。

早期抢占实验给四个 PCB 各编译进一块固定 20 KiB 存储，并在
低端留下一个 guard。这个实现没有地址申请、没有物理帧所有权，也没有真正的
销毁状态；容量和 PCB 槽位在链接时绑定。后来的实现保留四进程教学负载，却把栈
本身迁移成跨 KVA、buddy 和页表三层的动态对象。PCB 再拆为
Process/Thread 时，可以迁移所有者，而不必再次发明栈资源协议。

\topic{六页布局同时表达容量和两侧故障边界}
一个动态内核栈取得 6 个连续虚拟页，但只有中间 4 页存在叶映射：

\begin{center}
\begin{tikzpicture}[os diagram,node distance=0mm]
  \node[os hardware,minimum width=1.9cm] (lower) {lower guard\\4 KiB};
  \node[os state,right=of lower,minimum width=2.1cm] (d0) {data 0\\4 KiB};
  \node[os state,right=of d0,minimum width=2.1cm] (d1) {data 1\\4 KiB};
  \node[os state,right=of d1,minimum width=2.1cm] (d2) {data 2\\4 KiB};
  \node[os state,right=of d2,minimum width=2.1cm] (d3) {data 3\\4 KiB};
  \node[os hardware,right=of d3,minimum width=1.9cm] (upper) {upper guard\\4 KiB};
\end{tikzpicture}
\end{center}

设 KVA 返回的区间起点为 \(b\)，页大小 \(P=4096\)，则：
\[
\begin{aligned}
  \mathrm{lower\ guard} &= b,\\
  \mathrm{mapped\ begin} &= b+P,\\
  \mathrm{stack\ top} &= b+5P,\\
  \mathrm{upper\ guard} &= b+5P,\\
  \mathrm{range\ end} &= b+6P.
\end{aligned}
\]
可用容量为 \(4P=16384\) 字节，完整虚拟区间为 \(6P=24576\) 字节。x86-64
栈向低地址增长，所以 TSS.RSP0 指向 \(\mathrm{stack\ top}\)，处理器第一次
压栈落入 data 3。继续向下越过 data 0 会命中 lower guard。upper guard
不负责捕获正常的向下溢出；它负责隔开相邻对象，并让错误的正偏移、范围计算
或把“栈顶”误作可写首地址时立即得到 not-present 故障。

两个 guard 从申请到释放都属于同一 KVA 描述符，却始终没有 PTE。这正是
“not-present 不等于 free”的长期实例。若只保留 lower guard 而让下一对象
紧贴栈顶，错误的向上访问可能静默写进另一个内核对象；双 guard 用额外 4 KiB
虚拟地址换取对称且稳定的故障边界，并不消耗额外物理 RAM。

\topic{虚拟连续不要求四个物理帧连续}
\code{KernelStack} 保存一个六页
\code{KernelVirtualAddressRange}、四个 \code{PhysicalFrame}、已映射
页数和 active 状态。管理器对四个 data 页分别调用 order-0 物理分配：
\[
  V_{b+P+iP}\longmapsto F_i,\qquad i\in\{0,1,2,3\}.
\]
\(F_i\) 可以位于任意四个可用 PFN；页表把它们组织成连续的 16 KiB 虚拟栈。
若为每个小栈申请 order 2 连续块，长期运行后即使总空闲页很多，也可能因外部
碎片找不到四页连续物理区。order 0 更符合栈只要求虚拟连续的真实契约。

每帧在映射前通过 direct-map 清零。清零既避免新进程观察到旧内核栈残留，也
让首次构造现场以确定状态开始。销毁时再次清零后才归还 buddy，使后续所有者
不会继承寄存器、地址或用户参数。叶权限固定为 supervisor、RW、NX、
cache enabled；任何 user、executable 或 cache-disabled 偏差都由完整校验
拒绝。

管理器有 512 个长期槽，容量与 1024 项 KVA 描述符容量显式分开。槽索引由
Thread 持有，不等于 PID、TID 或数组身份；栈管理器本身不理解调度状态或
用户地址空间。统计记录：

\begin{osgridlongtable}{L{0.31\textwidth}L{0.53\textwidth}}
\textbf{字段} & \textbf{必须成立的含义}\\ \hline
\code{active\_stack\_count} & 当前 active 槽数量\\ \hline
\code{active\_mapped\_page\_count} & active 数量乘 4\\ \hline
\code{active\_guard\_page\_count} & active 数量乘 2，由前两者派生\\ \hline
\code{successful\_creation\_count} & 从初始化以来完整发布的栈数量\\ \hline
\code{destruction\_count} & 完整撤销并清空槽的数量\\ \hline
\code{peak\_active\_stack\_count} & 历史同时活动栈峰值\\ \hline
\code{peak\_active\_mapped\_page\_count} & 历史物理后备页峰值\\ \hline
\end{osgridlongtable}

累计计数必须满足：
\[
  \mathrm{creations}-\mathrm{destructions}
  =\mathrm{active\ stacks},
\]
\[
  \mathrm{active\ mapped\ pages}=4\times\mathrm{active\ stacks}.
\]
峰值只增不减，当前值不得超过峰值。64 位计数在提交前检查溢出，不能让一次
成功操作把守恒式绕回零。

\topic{创建是一笔直到发布前都可完整撤销的事务}
\code{TryCreate(slot)} 先验证管理器、槽位、依赖分配器和统计，再执行：

\begin{enumerate}[label=\arabic*.]
  \item 向 KVA 申请 6 页、1 页对齐的完整区间；
  \item 查询 6 个虚拟页，证明没有任何遗留叶映射；
  \item 逐页申请四个 order-0 物理帧并验证 direct-map 可访问；
  \item 清零当前帧，以 supervisor、RW、NX 权限映射对应 data 页；
  \item 核对两个 guard 仍 not-present、四个叶映射的物理身份和权限；
  \item 把候选对象复制到目标槽，最后设置 active 并更新统计。
\end{enumerate}

候选对象在最后一步之前不对调度器可见。第 \(j\) 个物理帧申请或映射失败时，
\code{RollbackCreation} 按相反顺序清除已经建立的叶项、清零并归还已经取得
的帧，最后精确释放六页 KVA。伪代码为：

\begin{verbatim}
range = kva.allocate(6)
for each data page:
    frame = buddy.allocate(order0)
    zero(frame)
    map(data_va, frame, supervisor | rw | nx)
validate(candidate)
publish(slot, candidate)

on any failure:
    unmap(mapped pages in reverse order)
    zero and release acquired frames in reverse order
    kva.release(exact range)
\end{verbatim}

回滚本身也可能失败，因此接口区分原始的
\code{FrameAllocationFailed} 或 \code{PageMappingFailed}，以及回滚时的
\code{RollbackFailed}。后一种状态说明资源集合已经无法证明，调用方不能把
它当成普通“内存不足”继续运行。失败输出不发布半有效 \code{KernelStack}，
目标槽仍保持全零。

\topic{六页内核栈为什么必须整段销毁}
仅验证六页都在 KVA 窗口内还不够。若其他错误代码提前释放了栈的 KVA
描述符，PTE 和物理帧可能仍然存在；此时销毁函数若先撤映射、归还物理帧，
最后才发现 KVA 不再归自己，就会留下无法重试的半销毁对象。

\code{KernelVirtualAddressAllocator::OwnsAllocation} 因而要求起点、页数和
Allocation 类型与活动描述符精确匹配。物理分配器的同名只读查询只接受精确
order-0 活动块；大块内部首帧或已释放帧都不满足。\code{ValidateStack} 在
任何破坏性步骤前同时验证两层所有权、两个 guard、四个叶映射、权限和四个
物理地址互异。只有预检完整通过，销毁才逆序：

\begin{verbatim}
unmap data pages from high to low
zero and release frames from high to low
release the exact six-page KVA allocation
clear the slot and commit statistics
\end{verbatim}

当前单 BSP、关中断提交边界让预检与撤销之间没有另一个管理器调用者。SMP
版本必须把“验证后执行”放入同一个锁或代际协议，否则预检成功不能阻止另一个
CPU 在提交前改变页表或 KVA 所有权。

\topic{高半区共享让所有进程在同一虚拟地址看见栈}
动态栈位于 KVA 窗口，对应 PML4 槽 402--465。创建进程根时，页表管理器复制
永久内核 PML4 的高半区项目；因此每个进程 CR3 都通过共享的内核页表子树看到
相同栈映射。栈不是用户地址空间中的私有叶页，也不需要为四个进程重复建立四份
映射。

共享带来两条严格顺序：

\begin{itemize}
  \item 创建栈必须发生在克隆进程根之前，使新根取得已经存在的高半区路径；
  \item 销毁用户地址空间只能释放低半区独占子树，不能递归释放 KVA 的共享
        高半区表。
\end{itemize}

当前高半区页表项本身由所有进程共享，所以在永久内核 CR3 下增加或清除叶项，
其他进程根会观察同一结构。单核路径通过 CR3 切换和 \code{INVLPG} 处理当前
TLB；SMP 后必须跟踪哪些 CPU 运行过相关根，并在物理帧复用前完成 shootdown。

\topic{176 字节首次现场必须完全落在最后一个 data 页}
\code{UserPrivilegeFrame} 长 176 字节。创建进程时，运行时读取管理器发布的
栈顶 \(t\)，按 16 字节 ABI 对齐后在其下方构造首次现场：
\[
  f=\operatorname{alignDown}(t-176,16).
\]
管理器的 \code{Contains(slot,f,176)} 证明半开区间
\([f,f+176)\) 完整落在四个 data 页中。当前布局使它落在最高 data 页，写入
物理内存时需使用该页对应的独立 frame 和页内偏移，不能假设四个 PFN 连续。

首次派发前，\code{ActivateProcess} 同时验证保存帧仍属于该槽、切换到进程
CR3，并把 TSS.RSP0 设置为该动态栈顶。抢占切换时返回的另一个帧地址和新的
RSP0 属于同一个目标进程；若只换 CR3 不换 RSP0，下次用户态陷入会把现场压到
旧进程的栈上。

\topic{退出后的栈什么时候才能释放}
用户 exit 或用户异常发生时，C++ 处理函数正在终止进程的动态栈上运行。
\code{CompleteCurrentProcess} 只完成以下工作：

\begin{enumerate}[label=\arabic*.]
  \item 保存终止现场并把调度状态改为 Terminated；
  \item 关闭描述符并选择后继；
  \item 切换到永久内核 CR3，销毁进程独占的用户页和页表；
  \item 有后继时切换其 CR3/RSP0 并返回其保存帧；无后继时恢复默认 RSP0；
  \item 由汇编恢复路径最终返回永久启动栈。
\end{enumerate}

此时终止栈仍映射且仍由管理器持有。只有
\code{OsKernelEnterScheduledProcess} 返回到永久启动栈之后，
\code{ReapTerminatedKernelStacks} 才扫描 Terminated 槽。回收器先读取真实
RSP，并用 \code{Contains} 检查一个 64 位探针范围；只要当前 RSP 仍在目标
栈内，就拒绝销毁。通过检查后才调用 \code{TryDestroy}，成功后把 PCB 中的
保存帧指针清为 null。

\begin{center}
\begin{tikzpicture}[os diagram]
  \node[os state] (user) {Ring 3\\运行};
  \node[os block,right=of user] (exit) {exit/\#UD/\#PF\\逻辑终止};
  \node[os state,right=of exit] (switch) {永久 CR3\\用户空间回收};
  \node[os block,below=of switch] (boot) {汇编返回\\永久启动栈};
  \node[os state,left=of boot] (reap) {RSP 安全检查\\动态栈销毁};
  \draw[os arrow] (user) -- (exit);
  \draw[os arrow] (exit) -- (switch);
  \draw[os arrow] (switch) -- (boot);
  \draw[os arrow] (boot) -- (reap);
\end{tikzpicture}
\end{center}

“当前 CR3 已经不是进程根”不能替代 RSP 检查，因为动态栈属于共享高半区，在
永久根里同样存在。“进程状态已经 Terminated”也不能替代检查，因为状态是
调度事实，不是 CPU 当前栈指针的物理事实。安全点必须由真实处理器状态证明。

\topic{完整校验从槽集合重建跨层事实}
\code{KernelStackManager::Validate} 不只比较计数。它遍历全部 512 槽并
重新证明：

\begin{enumerate}[label=\arabic*.]
  \item inactive 槽的 range、frame、页数和 active 位全部为零；
  \item active range 为 6 页、页对齐、无加法溢出且首尾均 canonical；
  \item KVA 分配器精确拥有这一个 allocation；
  \item lower/upper guard 的查询结果都是 NotMapped；
  \item 四个 data 页各自映射到记录的物理帧，页粒度为 4 KiB；
  \item 权限精确为 supervisor、RW、NX、普通缓存；
  \item 四帧均由 buddy 以 order 0 持有、direct-map 可访问且彼此不同；
  \item 重建的活动数、映射页数、累计差和峰值与缓存统计一致；
  \item 最后再运行 KVA 与 buddy 自身的完整不变量检查。
\end{enumerate}

跨层校验让一种资源的“局部合法”不能掩盖另一层漂移。例如一个 PTE 仍 present
但 KVA 描述符已丢失，在页表看来可以访问，在 buddy 看来帧也已分配，只有联合
所有权协议能判为损坏。

\topic{三层测试和三条整机路径怎样分工}
单元测试覆盖初始化失败、非法槽、重复创建/销毁、地址窗口耗尽、物理帧耗尽、
页表冲突、创建回滚、读出、范围包含、双 guard、权限、非连续物理帧、统计和
人为破坏。测试先外部释放一个物理帧，证明 PTE 仍在也不能冒充所有权；原址
重新申请后再外部释放 KVA，二者都要求 \code{Validate} 与销毁在破坏性操作
前拒绝。最后修复 KVA、篡改并恢复 PTE，只有三层重新一致后才允许安全销毁。

集成测试创建四个栈，逐一在最高 data 页构造 176 字节现场，并建立新的进程
页表根。它从该根再次查询所有 guard 与 data 映射，证明高半区共享不是单一
永久 CR3 下的偶然现象。逆序销毁后，四个物理帧全部清零，buddy 空闲数和
KVA 页数恢复空闲基线；共享页表根只保留一张仍可能被进程 PML4 引用的 PDPT，
临时 PT/PD 均已回收。创建、销毁均为 4，峰值活动栈为 4，峰值映射页为 16。

固定种子 \code{0x4B535441434B524E} 在 64 个槽和 4096 页独立所有权模型上
执行 100000 次随机创建/销毁。操作类型取随机数高位，槽索引使用下一次独立
随机值，避免线性同余发生器低位相关性伪造覆盖。参考模型自己寻找六页
best-fit 空洞；每 257 步比较映射、统计和三层不变量，最后随机状态全部排空。

QEMU 正常路径、用户 \#UD 隔离路径和用户 \#PF 隔离路径都必须经过同一动态栈
创建与回收协议。正常四进程路径要求每个 lower guard、栈顶和 upper guard
地址各出现一次；故障路径要求异常进程也回到安全点并释放栈。整机冷路径输出：

\begin{verbatim}
[OS][KERNEL] KERNEL_STACK_MANAGER_READY
[OS][KERNEL] KERNEL_STACK_SLOT_CAPACITY=0x0000000000000200
[OS][KERNEL] KERNEL_STACK_MAPPED_PAGES=0x0000000000000004
[OS][KERNEL] KERNEL_STACK_GUARD_PAGES=0x0000000000000002
[OS][KERNEL] KERNEL_STACK_SIZE_BYTES=0x0000000000004000
...
[OS][KERNEL] KERNEL_STACK_ACTIVE_STACKS=0x0000000000000000
[OS][KERNEL] KERNEL_STACK_SUCCESSFUL_CREATIONS=0x...
[OS][KERNEL] KERNEL_STACK_DESTRUCTIONS=0x...
[OS][KERNEL] KERNEL_STACK_PEAK_ACTIVE_STACKS=0x...
[OS][KERNEL] KERNEL_STACK_PEAK_MAPPED_PAGES=0x...
[OS][KERNEL] KERNEL_STACK_RESOURCES_RECLAIMED
\end{verbatim}

这些日志在准备阶段、每 Thread 创建结果和最终聚合点输出，不在页申请、映射、
时钟中断或回收循环中逐步输出。宿主 QEMU runner 继续给每行添加单调时间；
来宾使用 PIT 单调 tick 表达自己的时间域。丰富日志在这里意味着关键状态可
关联、计数可以核对，而不是把高频事件冲进 115200 波特串口。累计值随
4/4、64/128、256/512 三档 Process/Thread 容量事务变化，因此不再伪装成
单一常数；测试要求 successful creations 与 destructions 相等、active 为零，
并要求 peak mapped pages 等于 peak active stacks 的四倍。

\topic{当前完成边界与下一资源问题}
动态栈已经消除了编译期四份 20 KiB 进程栈存储。早期由四个 PCB 保存
管理器槽索引，存储上限为 256；现在所有者已经迁移为 Thread，槽上限
提升为 512。管理器仍是单 BSP、无内部锁；栈大小固定 16 KiB，不支持按需
增长，也没有每 CPU emergency stack 或 stack canary。空 PT/PD 与独占
PDPT 回收已经完成，并由页表帧统计证明动态栈反复创建/销毁不会永久增长
中间表基础设施。最后再把这些局部结果接入通用引用计数、
可组合的作用域回滚与跨层资源快照。
